\bookchapter{付録H：バージョンロードマップ}

\chapterintro{この本は更新を続けています。以下は今後のバージョンで計画している方向性です。フィードバックで「いちばん読みたい内容」を教えてもらえるとうれしいです。}

\textbf{v1.8 で実現：第7·8·9部『AIプログラミング実践編』『Android 開発実践編』『iOS 開発実践編』を新設（第37〜69章、全33章）}\par

AI プログラミングのパラダイム革命、5大思考原則、コンテキストと要件のエンジニアリング、2026年のツールとモデルの全景、API 実践、会話エンジニアリング、3つの完全なプロジェクト、品質保証とステップアップへの道——「AI を使える」から「AI に指揮をとってコードを書かせる」レベルへ引き上げます。さらに Android と iOS の2大プラットフォームの完全な開発チェーン：環境構築（Android Studio／Xcode）、言語速習（Kotlin／Swift）、UI フレームワーク（Jetpack Compose／SwiftUI）、データとネットワーク、実践プロジェクト（家計簿アプリ＋天気アプリ）、テストとリリース（Google Play／App Store）、そしてクイックリファレンス——デジタルユーザーから App 開発者へと一気に成長できます。

\begin{tcbinfo}{次回版（v1.9）の計画テーマ}
\begin{itemize}
  \item スマホとウェアラブル：Apple Watch、Huawei のスマートウォッチ、スマートリング（Oura／Galaxy Ring）を深く連携させるガイド
  \item 旅とスマホ：eSIM の実践、海外データ通信プラン、オフライン地図、海外での決済、翻訳 App の比較
  \item スマホと健康管理：運動と健康、睡眠モニタリング、心拍·血中酸素、女性の健康、メンタルヘルス App
  \item AI ショート動画づくり：剪映（CapCut）AI、Kling AI、デジタルヒューマン、AI ナレーションの実践チュートリアル
  \item 縦画面ショートドラマ入門：脚本の構成、撮影テクニック、各プラットフォームのポリシー
\end{itemize}

\end{tcbinfo}

\begin{tcbwarn}{情報の鮮度について}
スマホの技術は非常に速く進化するため、本書の内容は2026年半ばの技術と市場の状況に基づいて執筆しました。AI の能力、チップの型番、OS のバージョン、価格帯は、あなたが読む時点では変わっているかもしれません。具体的な操作手順については、お使いの端末の実際のシステムバージョンを優先してください。購入の判断には、最新のレビューとユーザーの声を確認しましょう。

\end{tcbwarn}
